
\documentclass[12pt, a4paper]{article}

\usepackage[margin=1in]{geometry}
\usepackage{setspace}
\usepackage{parskip}
\usepackage{booktabs}
\usepackage{array}
\usepackage{graphicx}
\usepackage{float}
\usepackage{adjustbox}
\usepackage{rotating}
\usepackage{caption}
\usepackage{subcaption}
\usepackage{amsmath}
\usepackage{amssymb}
\usepackage{hyperref}
\usepackage{natbib}
\usepackage{fancyhdr}
\usepackage{titlesec}
\usepackage{xcolor}
\usepackage{lscape}
\usepackage{microtype}

\onehalfspacing

\definecolor{accentblue}{RGB}{30, 80, 150}


\titleformat{\section}
  {\large\bfseries\color{accentblue}}
  {\thesection.}{0.75em}{}[\titlerule]

\titleformat{\subsection}
  {\normalsize\bfseries}
  {\thesubsection.}{0.75em}{}

\captionsetup{
  font=small,
  labelfont=bf,
  labelsep=period,
  justification=centering
}

\pagestyle{fancy}
\fancyhf{}
\fancyhead[L]{\small\textit{Medical Insurance Charges: Sex-Based Analysis}}
\fancyhead[R]{\small\textit{March 24, 2026}}
\fancyfoot[C]{\thepage}
\renewcommand{\headrulewidth}{0.4pt}

\hypersetup{
  colorlinks = true,
  linkcolor = accentblue,
  citecolor = accentblue,
  urlcolor = accentblue
}

\begin{document}

\begin{titlepage}
  \centering
  \vspace*{2cm}

  {\LARGE\bfseries Medical Insurance Charges:\\[0.4em]
  Sex-Based Analysis}

  \vspace{1.5cm}

	{\large Matthew deHaas, Callum Mackenzie, Saksham Saksham, Nihar Mahajan}

  \vspace{1.5cm}

  {\normalsize {Course: STAT 306}}\\[0.4em]
  {\normalsize Instructor: Kenny Chiu}

  \vfill

  {\normalsize March 24, 2026}

  \vspace{1cm}
\end{titlepage}

\section{Introduction}

This report investigates whether sex exerts predictive power on medical insurance premiums after controlling for baseline demographics, health metrics, and financial histories—specifically age, body mass index (BMI), risk score, and total claims paid. Motivation comes from a collective interest in whether non-meritorious, sex-based discrimination is present in the context of medical insurance applicants. Such an influence could potentially suggest that insurance institutions operate under systemically discriminatory practices as it pertains to sex. Data was observed from a collection of studies conducted by accredited institutions such as the World Health Organization, the National Institute of Health, the Centre for Disease Control and Prevention, and more. The dataset itself provides information on approximately 100,000 individuals, including their “...demographics, socioeconomic status, health conditions, lifestyle factors, insurance plans, and medical expenditures”. As mentioned, age and BMI will serve as a primary basis for analysis. It may be helpful to note that the source under scrutiny that collected these studies underwent a partial cleaning and normalization procedure, so certain pieces of data, particularly in the extremities, may seem strange in isolation. In addition, the data became much more coherent when beginning at the minimum value of the response variable, and performing a log transformation proved to be an important visual aid.

\section{Analysis}

\subsection{Univariate Summaries}

To begin, we introduce the variables used for analysis and their respective shapes.

\begin{figure}[H]
  \centering
  \includegraphics[width=0.75\textwidth]{images/insurance_charges_distribution.png}
  \caption{Distribution of annual insurance premiums. There is significant right skew, so perhaps a log transformation is more appropriate.}
  \label{fig:charges_dist}
\end{figure}

\begin{figure}[H]
  \centering
  \includegraphics[width=0.75\textwidth]{images/log_insurance_charges_distribution.png}
	\caption{Much more suitable for analysis.} 
  \label{fig:log_charges_dist}
\end{figure}

\begin{figure}[H]
  \centering
  \begin{subfigure}[b]{0.48\textwidth}
    \includegraphics[width=\textwidth]{images/age_distribution.png}
    \caption{Age}
    \label{fig:age_dist}
  \end{subfigure}
  \hfill
  \begin{subfigure}[b]{0.48\textwidth}
    \includegraphics[width=\textwidth]{images/bmi_distribution.png}
    \caption{BMI}
    \label{fig:bmi_dist}
  \end{subfigure}
  \caption{Distributions of BMI and Age. We observe that they both seem are roughly normal.}
  \label{fig:predictor_dists}
\end{figure}

\subsection{Bivariate Relationship}

We now explore how sex in isolation relates to insurance premiums.

\begin{figure}[H]
	\includegraphics[width=\textwidth]{images/log_insurance_charges_vs_sex_distribution.png}
	\caption{Log Insurance Premiums vs. Sex }
  \caption{Insurance premiums as a function of the sexes}
  \label{fig:predictor_dists}
\end{figure}

\subsection{Confounders}

This looks promising, but a simple male to female comparison is naive. We now consider how age and BMI are distributed across the sexes, as well as how they compare to (log) insurance premiums.

\begin{figure}[H]
	\centering

  \begin{subfigure}[b]{0.48\textwidth}
		\includegraphics[width=\textwidth]{images/age_vs_sex_distribution.png}
		\caption{Age vs. Sex}
		\label{fig:age_sex__dist}
  \end{subfigure}
  \hfill
  \begin{subfigure}[b]{0.48\textwidth}
		\includegraphics[width=\textwidth]{images/age_vs_log_charges_distribution.png}
		\caption{Age vs. Log Insurance Premiums}
		\label{fig:age_charges__dist}
  \end{subfigure}

\end{figure}

Observe that the distribution of age among the sexes is nearly identical as per the box plots. In addition, the relationship between age and age and premiums is weakly positive among the sexes. We can infer that age is an unlikely influence on the differences in premiums among the sexes.

\begin{figure}[H]
  \centering

  \begin{subfigure}[b]{0.48\textwidth}
		\includegraphics[width=\textwidth]{images/bmi_vs_sex_distribution.png}
		\caption{BMI vs. Sex}
		\label{fig:bmi_sex__dist}
  \end{subfigure}
  \hfill
  \begin{subfigure}[b]{0.48\textwidth}
		\includegraphics[width=\textwidth]{images/age_vs_log_charges_distribution.png}
		\caption{BMI vs. Log Insurance Premiums}
		\label{fig:bmi_charges__dist}
  \end{subfigure}

\end{figure}

Similarly, BMI itself and its relationship with premiums are weakly positive among the sexes. We can comfortably conclude now that the variables age and BMI under scrutiny here are not likely to be confounding in comparison between sex and medical insurance premiums. Now consider the risk score:

\begin{figure}[H]
  \centering
	\includegraphics[width=\textwidth]{images/risk_vs_logtotal.png}
	\caption{Risk score paid vs. Log Insurance Premiums}
	\label{fig:bmi_charges__dist}
\end{figure}

We can see that also demonstrates weakly positive correlation. Finally, the log total claims paid:

\begin{figure}[H]
  \centering
	\includegraphics[width=\textwidth]{images/logclaimspaid_vs_logtotal.png}
	\caption{Log Claims paid vs. Log Insurance Premiums}
	\label{fig:bmi_charges__dist}
\end{figure}

Visually this is the best fit we have encountered so far. There is a strong positive correlation in the data.



\subsection{The Case for Regression}

There are myriad factors to consider in attempting to understand the prices of medical insurance premiums, especially with the contentious notion of systemic discrimination. Thus, we look to linear regression to explain whether a sex is a primary determinant of what one is charged for medical insurance.

We will examine two models:
\begin{enumerate}
    \item $\log Y=\beta_0+\beta_1\cdot\text{Sex}$
    \item $\log Y=\beta_0+\beta_1\cdot\text{Sex}+\beta_2\cdot\text{Age}+\beta_3\cdot\text{BMI}+\beta_4\cdot \text{Risk Score}+\beta_5\cdot\log\text{Claims Paid}$
\end{enumerate}

Evidenced by our EDA plots, there is negligible interaction between sex and age or sex and BMI, so we are only considering models which have the sex covariate changing the intercept.

\newpage


% Table created by stargazer v.5.2.3 by Marek Hlavac, Social Policy Institute. E-mail: marek.hlavac at gmail.com
% Date and time: Sun, Mar 22, 2026 - 23:41:48
\begin{table}[!htbp] \centering 
  \caption{Staged Linear Regression Results} 
  \label{} 
\begin{tabular}{@{\extracolsep{5pt}}lcc} 
\\[-1.8ex]\hline 
\hline \\[-1.8ex] 
 & \multicolumn{2}{c}{\textit{Dependent variable:}} \\ 
\cline{2-3} 
\\[-1.8ex] & \multicolumn{2}{c}{log(annual\_premium)} \\ 
 & Sex Model & Sex + Age + BMI \\ 
\hline \\[-1.8ex] 
 Sex & 0.0046 & 0.0042 \\ 
  & (0.0031) & (0.0031) \\ 
  Age (Years) &  & 0.0048$^{***}$ \\ 
  &  & (0.0001) \\ 
  BMI (kg/m$^2$) &  & 0.0055$^{***}$ \\ 
  &  & (0.0003) \\ 
  Constant & 6.2260$^{***}$ & 5.8520$^{***}$ \\ 
  & (0.0022) & (0.0096) \\ 
 \hline \\[-1.8ex] 
Observations & 97,987 & 97,987 \\ 
R$^{2}$ & 0.00002 & 0.0278 \\ 
Adjusted R$^{2}$ & 0.00001 & 0.0278 \\ 
Residual Std. Error & 0.4846 (df = 97985) & 0.4778 (df = 97983) \\ 
F Statistic & 2.2422 (df = 1; 97985) & 935.4822$^{***}$ (df = 3; 97983) \\ 
\hline 
\hline \\[-1.8ex] 
\textit{Note:}  & \multicolumn{2}{r}{$^{*}$p$<$0.05; $^{**}$p$<$0.01; $^{***}$p$<$0.001} \\ 
\end{tabular} 
\end{table} 


In addition to the table above, we will note that the p-value for the intercept in the model with sex as the only covariate is $0.134$, and in the one with age, sex, and BMI, it is $0.164$.

We now have a clear picture of the data. First, we will check the hypothesis for the gender coefficient $H_0:\beta_1=0$, $H_1:B_1\neq0$. Given that the p-value for $\beta_0$ is below $0.05$ for both models, we accept the null hypothesis that there is no difference in the insurance premium between the sexes. This is consistent with the visual intuition obtained from the previous scatter plot that related age and BMI to log premiums separately. This is strengthened by the $R^2$, the variable responsible for the proportion of variance in log premiums explained by the covariates in the model. It was at a negligible 0.00002 in the sex model, while jumping to 0.1793 in the second one. 

We can see that other than sex, all of the other coefficients are statistically significant. Risk score and log claims paid have small nominal coefficients of $-0.0083$ and $0.0026$, respectively. Conversely, age and BMI have larger coefficients of $1.4854$ and $0.0428$, respectively. Under the hypothesis tests $H_0: \beta_4 = 0$, $H_1: \beta_4 \neq 0$, we reject the null hypothesis and state there is a statistically significant relationship between risk score and insurance premiums. Similar logic holds for log claims paid ($H_0: \beta_5 = 0$, $H_1: \beta_5 \neq 0$); we reject the null and conclude there is a relationship between historical claims paid and the resulting insurance premiums.

To evaluate our assumptions we will examine the Q-Q and residual-fitted value plots:

\begin{figure}[H]
  \centering

  \begin{subfigure}[b]{0.48\textwidth}
		\includegraphics[width=\textwidth]{images/qqplot.png}
		\caption{Q-Q Plot for Model 2}
		\label{fig:bmi_sex__dist}
  \end{subfigure}
  \hfill
  \begin{subfigure}[b]{0.48\textwidth}
		\includegraphics[width=\textwidth]{images/mrisk_residuals.png}
		\caption{Residual-Fitted Value Plot for Model 2}
		\label{fig:bmi_charges__dist}
  \end{subfigure}

\end{figure}

Our Q-Q plot does not have significant tails, so our normality assumption roughly holds. As well, our fitted value-residual plot does not have any obvious patterns, meaning our homoscedasicity assumption holds.

\section{Conclusion}

The motivating question for this report was whether sex was a significant predictor for medical insurance premiums when controlling for demographic conditions such as age and BMI. According to this study, we are confident in the assertion that there exists no strong relationship between what an individual is charged and their sex. However, the particular covariates used only account for under 18\% of the variance in insurance premiums, thus there may still exist unobserved confounding variables. This is a limitation of the our study. To address this, in future analysis, we should include additional variables in the model so to explain more of the data's variation such as smoking status, geographic region, previous health complications, since the core motivation is to isolate the unencumbered interaction between sex medical insurance premiums.
 However, we can state that there does exist a relationship between an individual's risk score and claims paid with their insurance premium.  

\section{Statement of AI}

Artificial Intelligence was leveraged to improve the style of the report, particular to assist in improving the visual quality and over all professionalism.

The following prompt was used: "How do i save a plot generated by qqnorm(residuals(mrisk), main = "Q-Q Plot for Model") qqline(residuals(mrisk), col="red") to disk?"

\section{Appendix}

\newpage
\bibliographystyle{plainnat}
\bibliography{references}

\end{document}
